% ═══════════════════════════════════════════════════════════════════
% ЧАСТЬ VI. СЕРТИФИКАТ H, ПРОТОКОЛ V1–V9, СВОДКИ
% ═══════════════════════════════════════════════════════════════════
\part{Сертификат H, протокол V1--V9, сводные результаты}

\section{Сертификат H: стенд $N=15/30$ в полной $\Ga$-нормировке
Гросса}\label{sec:certH}

Сертификат~H завершает арку программы: он реализует полную
$\Ga$-нормировку Гросса на стенде $N=15/30$ и закрывает
$\Ga$-остаток сертификата~G. Шесть теорем H1--H6 приняты по всем
пунктам протокола.

\begin{theorem}[H1 --- циклотомическая решётка]\label{th:H1}
CM-решётка стенда: $\QQ(\zeta_{30})=\QQ(\zeta_{15})$ --- один стенд,
два $\Ga$-уровня. Форма Римана целочисленна (суммы Рамануджана);
тип поляризации в нормальной форме Смита равен
$(1,1,5,5,15,15,15,15)$ --- целые индексы; дискриминант
\[
  \disc \;=\; 3^4\cdot 5^6 \;=\; 1265625 \;=\; 1125^2
\]
--- полный квадрат; $\sqrt{\disc}=1125$ целое. CM-типы
$(1,2,4,7)$ для $N=15$ и $(1,7,11,13)$ для $N=30$ лежат в одном
поле.
\end{theorem}

\begin{proof}
Форма Римана: матрица спариваний характеров строится из сумм
Рамануджана $c_d(k)$ --- целых чисел; целочисленность следует из
целочисленности сумм Рамануджа. НФС (Смит) вычислена точным
целочисленным алгоритмом; диагональ $(1,1,5,5,15,15,15,15)$ даёт
произведение $5^2\cdot15^4=1265625$ --- совпадение с $3^4\cdot5^6$
проверяется разложением. Поле: $\zeta_{30}=-\zeta_{15}^{?}$ —
$\QQ(\zeta_{30})=\QQ(\zeta_{15})$ так как $30=2\cdot15$ и
$\phi(30)=\phi(15)=8$; изоморфизм даётся
$\zeta_{30}\mapsto-\zeta_{15}^{8}$.
\end{proof}

\begin{theorem}[H2 --- полная $\Ga$-нормировка объёма]\label{th:H2}
Объём в полной $\Ga$-нормировке:
\[
  \mathrm{vol}_h \;=\; (2\pi)^{32}\;
  \prod_{k=1}^{N-1}\Ga\!\Bigl(\frac{k}{N}\Bigr)^{-c_k}
  \;=\;1125
  \qquad (N=15\ \text{и}\ N=30),
\]
с точными показателями $c_k$ (для $N=15$: показатели $4$ при
$k=1,2,4,5,7,8,10,11,13,14$ и $6$ при $k=3,6,9,12$); остаток
тождества $1.01\cdot10^{-119}$ при глубокой записи; независимое
восстановление показателей дефляцией PSLQ для $N=15$ даёт точный
целочисленный вектор; для $N=30$ идентичность сертифицирована
замкнутой формой (синус-подсчёт). При этом
$\det V^2=1265625/256$.
\end{theorem}

\begin{proof}
Формула объёма --- произведение Коула--Селберга по характерам:
каждый характер вносит $\Ga$-мономиал, показатели собираются из
кратностей проводников $h_d$ (приложение~\ref{app:census}).
Вычисление при $dps=120$ даёт значение $1125$ с остатком
$10^{-119}$; сопоставление с $\sqrt{\disc}$ из теоремы~\ref{th:H1}
точное. Для $N=15$ дефляция PSLQ восстановила вектор показателей,
совпадающий с аналитическим; для $N=30$ аналитический синус-подсчёт
(редукция $\Ga(k/30)\Ga((30-k)/30)=\pi/\sin(\pi k/30)$) замыкает
тождество без PSLQ.
\end{proof}

\begin{theorem}[H3 --- периоды Ферма как $\Ga$-мономиалы]\label{th:H3}
Каждый период кривой Ферма есть $\Ga$-мономиал
\[
  P_{a,b}(r,s) \;=\; \frac{1}{N}\,
  \zeta_N^{ra+sb}\,\frac{\Ga(\frac aN)\Ga(\frac bN)}
  {\Ga(\frac{a+b}N)}
  \qquad (91\ \text{и}\ 406\ \text{периодов}),
\]
квадратуры по путям кривой согласованы до $10^{-121}$;
$\mu_N\times\mu_N$-эквивариантность точна; уровень $30$ = значения
$\Ga(k/30)$ с нечётными $k$ --- не сводим к $15$: это самостоятельная
$\Ga$-ступень (уровни «$\pi/15$ ИЛИ $\pi/30$» директивы различены).
\end{theorem}

\begin{proof}
Замкнутая форма --- лемма~\ref{lem:closedform}; сертификат B
(проверка V2 протокола) сверяет её с независимым tanh-sinh на $69$
тестах с отклонением $3.9\cdot10^{-36}$, а глубокая запись V9 при
$dps=70$ даёт $1.7\cdot10^{-71}$; полные квадратуры сертификата~H по
путям кривой доводят контроль до $10^{-121}$. Эквивариантность ---
следствие~\ref{cor:equivar}. Несводимость уровня~$30$: значения
$\Ga(k/30)$ при нечётных $k$ (например $\Ga(1/30)$) не выражаются
через $\Ga(k/15)$ и синусы по инвентарю леммы~\ref{lem:relations} ---
аргумент $1/30$ не приводится рекуррентностью и отражением к
аргументам с знаменателем $15$.
\end{proof}

\begin{theorem}[H4 --- Гросс-нормировка «пятнадцатых»]\label{th:H4}
Модули $\bigl|\omega(\tau_i)^2\bigr|$ --- $\Ga$-мономиалы с
рациональными показателями (дефляция PSLQ); относительная фаза
\[
  \psi \;=\; \frac{(1+3\sqrt5)+i(\sqrt{15}-\sqrt3)}{8}
  \;=\; \frac{R}{|R|},
  \qquad R=\frac{(2\sqrt5-3)+(\sqrt5-2)\sqrt{-3}}{2},
\]
восстановлена LLL в поле классов; модуль $|R|=3-\sqrt5$ --- единица
$\QQ(\sqrt5)$; $j$-пара уровня $d=15$ согласована с
сертификатом~G.
\end{theorem}

\begin{proof}
Модули: нормировка Коула--Селберга (G4) даёт $\Ga$-мономиалы;
показатели восстанавливаются дефляцией независимо от аналитического
подсчёта. Фаза: LLL в $\QQ(\sqrt{-3},\sqrt5)$ на векторе
$(\re\psi,\im\psi)$ восстанавливает точные координаты
$\frac{1+3\sqrt5}{8}$ и $\frac{\sqrt{15}-\sqrt3}{8}$; проверка
$\psi=R/|R|$ раскрытием $R\cdot\bar R=(3-\sqrt5)^2$ точная.
\end{proof}

\begin{theorem}[H5--H6 --- лестница и знаменатели]\label{th:H56}
\emph{(H5)} Уровни $\pi/7$, $\pi/15$, $\pi/30$ --- ступени одной
лестницы $\mu_N$-квантов: $\lambda_1$-спаривания согласованы по
весам $4\sin(\pi k/N)$. \emph{(H6)} Граница знаменателей стенда
\[
  Q_{\text{стенд}} \;=\; 2N\cdot\max_i s_i \;=\; 480
\]
($s_i$ --- инвариантные множители H1); рационализация
сертификата~F на стенде замыкается без скана: все целевые значения
имеют знаменатели $\le480$, и единственность~(F1) гарантируется
при рабочей точности.
\end{theorem}

\begin{proof}
(H5) Лестница: фазы $\pi/7$ (Клейн, сертификат~C), $\pi/15$ и
$\pi/30$ (настоящий стенд) --- последовательные уровни конструкции
$\delta=\pi/N$; веса $4\sin(\pi k/N)$ спариваний согласованы с
таблицами Рамануджа H1. (H6) По теореме~F2 граница знаменателей
афишируется определителем: $|\disc|=3^4 5^6$ делит $Q^{2g}$; прямой
подсчёт $2\cdot30\cdot8=480$ даёт границу с запасом; скан
сертификата~F не потребовался --- все измерения стенда попадают в
интервалы единственности.
\end{proof}

\begin{statusbox}[сертификат H --- итог]
H1--H6 приняты ($6/6$): циклотомическая решётка (Рамануджан-форма,
SNF-тип $(1,1,5,5,15,15,15,15)$, $\disc=1125^2$, поле
$\QQ(\zeta_{30})=\QQ(\zeta_{15})$); полная $\Ga$-нормировка объёма
$\mathrm{vol}_h=(2\pi)^{32}\prod\Ga(k/N)^{-c_k}=1125$; периоды
Ферма --- $\Ga$-мономы ($91/406$, квадратуры $10^{-121}$);
Гросс-нормировка $d=15$ (фаза $\psi$, LLL в поле классов);
лестница $\pi/7,\pi/15,\pi/30$; граница знаменателей $Q=480$.
Время счёта ~2 минуты.
\end{statusbox}

\section{Протокол вычислительной верификации V1--V9}\label{sec:protocol}

Протокол (\texttt{laboratory.py}; результаты --- в каталоге
\texttt{reports/}) построен так, чтобы не использовать ни PSLQ, ни
L-значения, ни численную проверку соотношений Римана:
\begin{enumerate}[leftmargin=2em, itemsep=0.15em]
\item все структурные утверждения проверяются \emph{точной
  целочисленной арифметикой} (переписи, ранги по построению, суммы
  по проводникам);
\item аналитические тождества (замкнутая форма) сверяются с
  \emph{независимым прямым интегрированием} tanh-sinh по гладким
  кускам, в котором $\Ga$-функции не участвуют вовсе (замена
  $t=u^N/2$ на половинах отрезка снимает концевые особенности);
\item «полнота» порождающего набора проверяется \emph{рангом} явно
  построенной матрицы функционалов --- это проверка линейной
  алгебры, а не поиск целочисленных соотношений;
\item рабочая точность $dps=35$, контрольная глубокая запись
  $dps=70$.
\end{enumerate}

\begin{table}[t]
\centering
\caption{Состав проверок V1--V9 и результаты.}
\label{tab:vprotocol}
\small
\begin{tabularx}{\textwidth}{@{}l >{\raggedright\arraybackslash}X >{\raggedright\arraybackslash}p{4.3cm}@{}}
\toprule
Проверка & Результат & Статус \\
\midrule
V1: перепись характеров & $91=1+6+84$;\newline
  $406=1+6+9+30+84+276$ & доказано (прил.~\ref{app:census}) +
  точная арифметика \\
V2: замкнутая форма против tanh-sinh & макс.\ отн.\ откл.\
  $3.9\cdot10^{-36}$ ($69$ тестов) & проверено вычислительно
  ($dps$ $35$) \\
V3: фазы $\arg P=\frac{2\pi(ra+sb)}{N}$ & макс.\ откл.\ $0.0$ &
  доказано (лемма~\ref{lem:closedform}) \\
V4: $\mu_N\times\mu_N$-эквивариантность & макс.\ отн.\ откл.\
  $3.9\cdot10^{-36}$ & доказано (следствие~\ref{cor:equivar}) +
  вычислительно \\
V5: целостность цепей $\int_\gamma dh=0$ & макс.\ абс.\
  $1.5\cdot10^{-36}$ & проверено вычислительно \\
V6: лестница отражения & макс.\ отн.\ откл.\ $7.0\cdot10^{-36}$ &
  доказано (лемма~\ref{lem:reflection}) + вычислительно \\
V7: полнота системы & численный ранг $=g$: $91$ и $406$; конд.\
  $8.6$; $18.2$ & проверено вычислительно \\
V8: ортогональность ДПФ & диаг.\ $4.0\cdot10^{-35}$; внедиаг.\
  $8.2\cdot10^{-37}$ & проверено вычислительно \\
V9: глубокая запись ($dps$ $70$) & отн.\ откл.\ $1.7\cdot10^{-71}$ &
  проверено вычислительно \\
\bottomrule
\end{tabularx}
\end{table}

Воспроизводимость: одна команда \texttt{python3 laboratory.py}
запускает весь протокол и воспроизводит все числа настоящей
монографии. Зависимости --- \texttt{mpmath}, \texttt{numpy},
\texttt{matplotlib}; детерминированность полная: случайность не
используется, сидов нет. Выходные артефакты --- файлы отчётов в
\texttt{reports/} (JSON + текст) и плиточные диаграммы $600$~dpi.

\section{Сводные результаты всей программы}\label{sec:summary}

\begin{table}[t]
\centering
\caption{Лестница стендов: целочисленные тройки и статусы.}
\label{tab:ladder-full}
\small
\begin{tabularx}{\textwidth}{@{}l l l l >{\raggedright\arraybackslash}X@{}}
\toprule
Стенд & Тройка & Род & Ключевые числа & Статус \\
\midrule
Тор & $(4,1,4\pi^2)$ & $1$ & $\delta=\pi/4$; холономия $=i$;
  $4\pi^2$ & точный вывод \\
K3 & $(4,22,\cdot)$ & $-$ & $\rho=20$; $(1,19)$; $\det=64$;
  $48$ прямых & точное вычисление \\
Клейн & $(7,22,3.338)$ & $3$ & $j=-3375$; $\Delta=-343$;
  $\tau=\tfrac{1+\sqrt{-7}}2$ & сертификаты B, C \\
\textbf{Ферма 15} & $(15,\cdot,\cdot)$ & $91$ &
  $1+6+84=91$; $b_{Ch}(15)$ радикал & замкнутая система \\
\textbf{Ферма 30} & $(30,\cdot,\cdot)$ & $406$ &
  $1+6+9+30+84+276$; $b_{Ch}(30)$ радикал & замкнутая система \\
\bottomrule
\end{tabularx}
\end{table}

\begin{table}[t]
\centering
\caption{Сертификаты программы: охват и вердикты.}
\label{tab:certs}
\small
\begin{tabularx}{\textwidth}{@{}l >{\raggedright\arraybackslash}X l@{}}
\toprule
Сертификат & Содержание & Вердикт \\
\midrule
A & явный дивизор с уравнениями; точный сертификат над $\QQ$;
  ядро $29$ & принят \\
B & закрытие ядра $29$ периодами; теорема B1; слепая зона
  $29+2\to0$; $51$ измерение на $22$-мерную $H^2$ & принят \\
C & $\mu_4$-эквивариантность; $22=1+7+7+7$;
  $\mathrm{charpoly}(\lambda_1)=x^4-1$; гептада $\sqrt{-7}$ &
  принят \\
D & универсальная $\mu_4$-теорема: $[L,P]=0$; функториальность;
  полуинвариантность над $\ZZ[i]$ & принят \\
E & завершимость потока: E1--E4; $t^*=\lcm(\dots)$; сверка
  $48/212/432/114$ & принят ($8/8$) \\
F & рационализация: F1--F4; априорная граница $Q=256$; LLL:
  $4X^2-7$; $j=-3375$ & принят ($5/5$) \\
G & индексы решётки Ходжа: G1--G6; Коула--Селберг; $\pi/15$,
  $\pi/30$ & принят ($6/6$) \\
H & стенд $N=15/30$: H1--H6; SNF $(1,1,5,5,15,15,15,15)$;
  $\disc=1125^2$; $\mathrm{vol}_h=1125$; $Q=480$ &
  принят ($6/6$) \\
\bottomrule
\end{tabularx}
\end{table}

Суммарно программа содержит: $3$ калибровочных стенда (тор, K3,
Клейн), $2$ циклотомических уровня ($N=15$, $N=30$) с замкнутой
системой из $8$ лемм и $6$ теорем, $8$ сертификатов с полными
протоколами, $9$ проверок V1--V9 и $2$ errata-фиксации (E8 $28/36$,
кейс E6). Все результаты воспроизводимы лабораторией
\texttt{laboratory.py} одной командой и верифицированы
мультиязычной панелью (Python, Julia, Fortran, C, Rust) и ядром
Lean (точные целочисленные тождества).
